\scanpage{124}

La commutativité de $Z'$ finalement résulte déjà des axiomes $T_0\,3,T_0\,4$ ($T_0\,1$ et \uline{$T_0\,2$} sont \unsure{inutiles} !), pour un 2-polyèdre triangulé connexe, \deleted{\unsure{dans|dont}} \unsure{le cas} \uline{où} un sommet n'est pas déterminé par son étoile, alors le polyèdre est une double pyramide \inserted{(à base un $n$-gone ($n\geqslant3$))} \unsure{pour $n\ne4$ son} dont le groupe des automorphismes est \deleted{\unsure{le}}\inserted{pr.duit} du groupe diédral $\mathbb D_n$ \inserted{($\simeq G_0^+$)} par le groupe $\{1,\underline-\}$ engendré par l'antipodisme.

% slanted
\margin{sinon pour \unsure{déjà} introduire $Z'$ comme \unsure{commutant} de $G_0^+$ -- \unsure{alors qu'on aurait pu} le définir comme \unsure{celui} de $G_0$ tout \unsure{entier}\ldots}

Le commutant de $G_0^+$ \inserted{dans $\unsure{\mathfrak S_S|G_S}$} est réduit à : \deleted{$\unsure{\{1,\sigma\}|\{1,\varepsilon\}}\times\{1,\underline-\}$ pour $n$ impair (centre de} \inserted{$\mathfrak z\times\{1,\underline-\}$, où $\mathfrak z$ est le} centre de $G_0^+\simeq\mathbb D_n$ (réduit à $\{1\}$ si $n$ impair, et à $\mathbb Z/2\mathbb Z$ si $n$ pair) donc dans ce cas le commut\textsuperscript{ant}\ \unsure{=}\ est commutatif.

% slanted
\margin{le commutant de $G_0^+$ est \inserted{bien} égal à celui de $G_0$, et égal au centre de $G_0$\unsure{,}}

Reste le cas (plus intéressant pour nous) où $n=4$, i.e. celui où tous les sommets ont = ordre (c'est le cas de l'octaèdre. Dans \unsure{le|le cas} le commutant de $G_0^+$ \inserted{$\simeq\mathfrak S_4$} (égal du reste à celui de $G_0\simeq G_0^+$, et au centre de $G_0$) est réduit à $\{1,\underline-\}$. \uline{En tous cas $Z'$ est commutatif}.

% slanted
\margin{En tous cas, le commutant de $G_0^+$ \deleted{est égal au} centre de $G_0$, est, soit égal à $\{1,\underline-\}$ (cas $n$ impair ou $n=4$) soit à $D_2=\{1,\sigma\}\times\{1,\underline-\}$ (cas $n$ pair $\geqslant6$) -- \unsure{à ch.} un d'élément du commutant étant formé d'éléments de $G_0$ involutifs, satisfait trivialement $T_0\,3$. Mais dans le cas $n$ pair $\geqslant6$, la structure de graphe $A_\sigma$ déduite de $A_0$ en appliquant $\sigma$ (fixant les sommets $s,s'$ de la double pyramide, et transformant chaque sommet de $E$ en son opposé) ne correspond \unsure{,|qu'}à une structure \unsure{triangulée|de} \deleted{\unsure{de}} que si $n\equiv0\ (4)$}

Donc ds le cas symétrique [où tous les pts ont = ordre], on voit

\[
\left\{\begin{aligned}
a)&\ \text{ou bien tt \deleted{point}\inserted{sommet} est déterminé par son étoile,}\\
b)&\ \text{ou bien }P_0\text{ est un octaèdre, }\inserted{\text{(cas régulier)}}\
 Z'=\{1,\underline-\},\quad\underline-\\
 &\text{l'antipodisme, i.e. l'unique élément}\\
 &\text{central }\ne\mathrm{Id}\text{ de }G_0=\operatorname{Aut}(P_0)
 (\simeq\mathfrak S_4\times\{\pm1\}).
\end{aligned}\right.
\tag*{$(49)$}
\]

Eh bien on suppose \deleted{\unsure{pas}} que $T_0\,3,T_0\,4$, (pas que tous les sommets de = ordre), alors $P_0$ est la double pyramide, $Z'=\{1,\underline-\}$ si $n$ impair ou $n=4$, $Z'\simeq\{1,\sigma\}\times\{1,\underline-\}$ dans le cas $n$ pair $\geqslant5$ $]\!]$
% archive sheet number: 123
